\documentclass[conference]{IEEEtran}
\IEEEoverridecommandlockouts
\usepackage{amsmath,amssymb}
\usepackage{booktabs}
\usepackage{array}
\usepackage{xcolor}
\usepackage[hidelinks]{hyperref}

\begin{document}

\title{Retrieval-Augmented Question Answering over Heterogeneous, Multi-Source Knowledgebases: A Survey}

\author{
\IEEEauthorblockN{Srujan Zanjal}
\IEEEauthorblockA{Independent Project \\
Email: srujan.zanjal@infocepts.com}
}

\maketitle

\begin{abstract}
Retrieval-augmented generation (RAG) grounds a language model's output in text retrieved from an external corpus, and has been studied extensively for the case of a single, roughly homogeneous document collection. Real knowledgebases are rarely that uniform: websites, PDFs, video transcripts, and code repositories differ in how they are extracted, how they should be chunked, what a citation to them looks like, and what kind of question is likely to be asked of them. This survey organizes the techniques that make grounded question answering possible across such heterogeneous, multi-source knowledgebases: dense retrieval and RAG foundations, hybrid lexical-dense reranking, chunking strategies for long and structured documents, modality-specific retrieval challenges for web documents, video/transcript QA, code-repository retrieval, and multi-document fusion, and confidence estimation and abstention in retrieval-grounded generation. We use SAGE (Semantic Analysis and Generation Engine), a system that ingests and answers questions over crawled websites, PDF documents, YouTube video transcripts, and GitHub repositories through one shared pipeline, as a single illustrative case study near the end of the survey, showing where its concrete engineering choices --- a deterministic overview-chunk mechanism and a diversity floor over reranked results --- instantiate ideas discussed earlier in the survey rather than departing from them. We close with open challenges in cross-modal retrieval fusion, confidence calibration, and incremental re-indexing that remain unresolved across the systems and papers surveyed.
\end{abstract}

\begin{IEEEkeywords}
retrieval-augmented generation, multi-source retrieval, multi-modal question answering, survey, reranking, chunking
\end{IEEEkeywords}

\section{Introduction}

A retrieval-augmented generation (RAG) system answers a question by first retrieving passages from an external corpus and then conditioning a generator on them, rather than relying only on knowledge encoded in the generator's parameters \cite{lewis2020rag}. The corpus in most RAG research and most production deployments is a single, relatively uniform collection: a Wikipedia dump, a set of internal wiki pages, or a company's PDF manuals. This uniformity is convenient --- one chunking strategy, one citation format, one notion of ``a document'' --- but it is not how people actually accumulate the things they want to ask questions about. A student researching a topic has a lecture recording, a handful of assigned readings, and a related open-source project open at the same time; a new engineer onboarding onto a codebase has the project's marketing site, its documentation, a walkthrough video, and the repository itself. Answering one conversational question against whichever of these sources actually contains the relevant evidence, with a citation that correctly identifies which source and where, requires retrieval and generation techniques that were mostly designed and evaluated one modality at a time to be composed into a single, coherent pipeline.

This survey organizes the pieces of that composition. We cover, in turn, the retrieval and generation foundations that any RAG system builds on (Section~\ref{sec:foundations}); reranking and hybrid lexical-dense retrieval, which becomes more important, not less, once a retrieved candidate pool may be drawn from structurally different sources with different lexical characteristics (Section~\ref{sec:reranking}); chunking strategies, where the right unit of retrieval for a code file, a video transcript, and a prose document are not obviously the same thing (Section~\ref{sec:chunking}); and the modality-specific retrieval literatures for web documents, video and transcript QA, code-repository-aware retrieval, and multi-document fusion across heterogeneous sources (Section~\ref{sec:modality}). We then cover confidence estimation and abstention (Section~\ref{sec:confidence}), which matters more once a system must decide not just \emph{what} to retrieve but whether it retrieved \emph{enough} to answer responsibly across sources whose retrieval quality may differ. Section~\ref{sec:casestudy} uses one system, SAGE (Semantic Analysis and Generation Engine), as a single illustrative case study of how these pieces compose in practice --- SAGE ingests crawled websites, PDF documents, YouTube video transcripts, and GitHub code repositories into a common pipeline, and we use its two most concrete design decisions, a deterministic source-overview chunk mechanism and a diversity floor over reranked results, as worked examples of ideas introduced earlier in the survey, not as the survey's subject. Section~\ref{sec:challenges} discusses open problems that recur across the systems and papers we survey, and Section~\ref{sec:conclusion} concludes.

\textbf{Scope and relation to prior surveys.} General RAG design --- what to retrieve, when, how many passages, how to fuse and iterate --- is already surveyed at length \cite{gao2023ragsurvey}, and hallucination in generation more broadly has its own dedicated survey \cite{ji2023hallucination}; we do not attempt to duplicate either. This survey's angle is narrower and, we believe, complementary: rather than surveying RAG technique-by-technique across an assumed-uniform corpus, we organize technique choices around the specific seams that appear once a system's corpus is actually several structurally different corpora combined at query time --- chunking, reranking, and confidence calibration, in particular, all needed to be reconsidered in Sections~\ref{sec:reranking}--\ref{sec:confidence} precisely because the single-corpus assumption those seams are usually described under no longer holds. We selected the literature discussed below by starting from each of the four modality-specific retrieval literatures in Section~\ref{sec:modality} (web, video/transcript, code, multi-document) and working backward to the shared retrieval, reranking, chunking, and confidence machinery each of them assumes, rather than starting from one venue's or one year's proceedings; we make no claim of exhaustive coverage of any single one of these four literatures on its own terms.

\section{Problem Definition and Terminology}
\label{sec:problem}

Before surveying techniques, it is worth stating precisely what changes when a RAG system moves from one corpus to several heterogeneous ones, since much of the terminology in the RAG literature implicitly assumes a single corpus and does not distinguish these concerns.

We use \emph{source} to mean one ingested unit with its own identity and extraction process --- one crawled website, one PDF, one video, one code repository --- and \emph{knowledgebase} to mean a set of sources a user has grouped together. A \emph{chunk} is the unit actually embedded and retrieved, produced from a source by some modality-specific chunking function. A \emph{query} may target one source, an entire knowledgebase, or an explicit, user-selected subset of sources spanning one or more knowledgebases; we call the last case a \emph{combined-source query}, and it is the case that introduces requirements a single-corpus system does not have.

Three properties distinguish multi-source retrieval from ordinary multi-document retrieval within one corpus:

\begin{itemize}
\item \textbf{Heterogeneous extraction.} Each source type has its own extraction pipeline (Section~\ref{sec:modality}) producing text of different lexical character --- crawled prose, PDF-extracted prose with layout artifacts, ASR-transcribed speech, and source code --- before any retrieval model sees it.
\item \textbf{Heterogeneous chunking.} As Section~\ref{sec:chunking} discusses, the natural retrieval unit differs by modality (a paragraph, a timestamped transcript segment, a function-scoped code block), so a single global chunking function is a simplification, not a given.
\item \textbf{Source-level accountability.} A citation in a multi-source system must identify which source, not just which passage, produced a piece of evidence, and a combined-source query carries an implicit expectation that every source the user explicitly selected was actually considered --- a requirement with no analogue when ``the corpus'' is one undifferentiated collection.
\end{itemize}

Table~\ref{tab:overview} previews how the sections that follow map onto these three properties, and we refer back to it as a running structure for the rest of the survey.

\begin{table}[t]
\caption{How survey sections address the three distinguishing properties of multi-source retrieval}
\label{tab:overview}
\centering
\scriptsize
\begin{tabular}{@{}>{\raggedright\arraybackslash}p{1.5cm}>{\raggedright\arraybackslash}p{2.6cm}>{\raggedright\arraybackslash}p{3.6cm}@{}}
\toprule
Section & Property addressed & Representative technique \\
\midrule
\ref{sec:foundations} & foundation shared by all sources & dense bi-encoder retrieval \cite{karpukhin2020dpr}, RAG \cite{lewis2020rag} \\
\ref{sec:reranking} & source-level accountability & rank fusion \cite{cormack2009rrf}, hybrid scoring \cite{gao2020clear} \\
\ref{sec:chunking} & heterogeneous chunking & late chunking \cite{gunther2024latechunking}, per-domain sizing \cite{jimenoyepes2024chunking} \\
\ref{sec:modality} & heterogeneous extraction & per-modality retrieval literatures \\
\ref{sec:confidence} & source-level accountability & selective prediction \cite{kamath2020selective} \\
\bottomrule
\end{tabular}
\end{table}

\section{Dense Retrieval and RAG Foundations}
\label{sec:foundations}

Neural information retrieval replaced sparse lexical matching with learned dense representations of queries and documents, beginning with semantic-matching architectures such as DSSM \cite{huang2013dssm} and maturing into dual-encoder retrievers such as Dense Passage Retrieval (DPR) \cite{karpukhin2020dpr}, which encodes queries and passages independently into a shared embedding space and retrieves by inner product or cosine similarity. This bi-encoder design is what makes retrieval tractable at corpus scale: passage embeddings are computed once, offline, and indexed for approximate nearest-neighbor search, so query-time cost does not grow with corpus size the way a query-passage joint encoder's would.

Retrieval-augmented generation \cite{lewis2020rag} formalized combining such a retriever with a generator trained (or prompted) to condition its output on the retrieved passages, marginalizing over retrieved evidence rather than depending solely on parametric memory. Fusion-in-Decoder \cite{izacard2021fid} showed that a sequence-to-sequence generator can encode many retrieved passages independently and let the decoder attend across all of them jointly, which scales better with the number of retrieved passages than concatenating everything into one encoder pass. The broad design space that has grown around this core idea --- what to retrieve, when to retrieve, how many passages, how to fuse them, whether to iterate retrieval and generation --- is surveyed at length elsewhere \cite{gao2023ragsurvey}; we do not repeat that survey's scope here, but note that essentially every choice it catalogs (retrieval granularity, fusion strategy, iteration) must be made \emph{per modality} once a system ingests more than one kind of source, which is the throughline of this paper.

Underlying any dense retriever is a sentence or passage embedding model. Sentence-BERT \cite{reimers2019sbert} established that fine-tuning a Siamese BERT encoder with a similarity objective produces embeddings whose cosine distance is directly meaningful for retrieval and clustering, without the impractical quadratic cost of scoring every query-passage pair with a full cross-encoder. Distilled variants such as MiniLM \cite{wang2020minilm,wang2020minilmv2} compress a large teacher model's attention behavior into a much smaller student, trading a modest amount of accuracy for substantially lower inference latency and memory footprint --- a tradeoff that matters directly for any system where retrieval must complete within an interactive, conversational time budget rather than an offline batch job. The specific distilled sentence-encoder checkpoint \texttt{all-MiniLM-L6-v2} \cite{minilm_modelcard} is a widely deployed example of this class of model in open-source RAG tooling.

Two further design choices in this foundational layer matter more, not less, once a system indexes several source types together. First, the choice of approximate nearest-neighbor index (commonly HNSW-graph-based, as in most production vector databases) is normally treated as an implementation detail orthogonal to retrieval quality, but its distance metric --- cosine versus inner product versus Euclidean --- must be fixed once for the entire index if all sources share one collection, which implicitly assumes that a sentence encoder's embedding geometry behaves comparably whether the encoded text came from prose, transcribed speech, or source code; this assumption is rarely tested directly in the systems we are aware of, our own case study in Section~\ref{sec:casestudy} included. Second, DPR's original training recipe \cite{karpukhin2020dpr} relies on in-batch and hard negative sampling drawn from one corpus's question-passage pairs; a multi-source system that reuses an off-the-shelf sentence encoder rather than fine-tuning per deployment (as our case study does) inherits whatever negative-sampling distribution the encoder's original training data reflected, which is unlikely to match the true mixture of prose, transcript, and code negatives such a system will actually encounter at query time. We are not aware of published work that fine-tunes a single encoder against a deliberately multi-modal negative distribution of this kind, which we return to as an open problem in Section~\ref{sec:challenges}.

\section{Reranking and Hybrid Retrieval}
\label{sec:reranking}

Dense retrieval alone systematically underperforms on queries dominated by exact terms --- identifiers, proper nouns, error codes, function names --- that a bi-encoder's semantic similarity does not weight as heavily as a human would. Two broad families of technique address this. The first re-scores an initial dense-retrieved candidate list with a more expensive, jointly-encoded model: cross-encoder rerankers \cite{nogueira2019bert} score a query-passage pair together rather than independently, and late-interaction models such as ColBERT \cite{khattab2020colbert} compute token-level similarity between query and passage representations, retaining much of a cross-encoder's accuracy at lower cost by deferring interaction to a lighter final step. The second family blends dense and lexical signals directly, either as a learned combination such as CLEAR \cite{gao2020clear}, which trains a dense retriever to model exactly the residual that a classical lexical scorer such as BM25 \cite{robertson2009bm25} fails to capture, or as an unweighted fusion of independently produced rankings, as in Reciprocal Rank Fusion \cite{cormack2009rrf}, which combines ranked lists from different retrievers or fields using only rank position rather than requiring calibrated, comparable scores. Concretely, RRF scores a document $d$ as
\[
\mathrm{RRF}(d) = \sum_{i} \frac{1}{k + \mathrm{rank}_i(d)},
\]
summed over each ranked list $i$ in which $d$ appears, with $\mathrm{rank}_i(d)$ its position in list $i$ and $k$ a small constant (commonly 60) that damps the influence of any single list's very top ranks; a document absent from a given list simply contributes zero from that list rather than requiring an imputed score. This rank-only property is exactly what makes RRF attractive for fusing across sources with different lexical and semantic characteristics (Section~\ref{sec:modality}): it needs no assumption that a 0.73 semantic score from a prose retriever means the same thing as a 0.73 score from a code retriever, only that each retriever's own internal ranking is locally meaningful.

This second family is a natural fit for multi-source retrieval specifically, for two reasons a single-corpus system does not have to confront. First, when a query's candidate pool is drawn from several structurally different sources at once --- prose paragraphs, code fragments, transcript segments --- their embeddings and lexical statistics were produced by extraction and chunking processes with different characteristics (Section~\ref{sec:chunking}), so a rank-based or explicitly weighted fusion is more robust to those characteristics differing across sources than a score that assumes one uniform passage distribution. Second, once ranking output must additionally decide not just \emph{which} passages are best but \emph{whether at least one passage from each explicitly requested source} should be represented in the final answer, ranking becomes a two-stage problem: score, then select subject to a source-coverage constraint, rather than score-and-truncate alone. We return to a concrete instance of this in Section~\ref{sec:casestudy}.

It is worth being explicit about the cost asymmetry between these two families, since it drives which one a given system can afford. A cross-encoder or ColBERT-style pass \cite{nogueira2019bert,khattab2020colbert} requires a forward pass through a neural network for every candidate under consideration, which is straightforward to justify when reranking, say, the top 100 candidates from a single retriever once per query, but becomes proportionally more expensive exactly when a multi-source system widens its candidate pool to guarantee per-source representation (Section~\ref{sec:casestudy} gives a concrete pool-widening formula from our case study): a wider pool for fairness purposes directly inflates the cost of any reranking stage that scales per-candidate. A closed-form hybrid score, by contrast, is computed from quantities the retrieval step already produced (a similarity score) plus a cheap token-set operation, so widening the candidate pool for source-fairness reasons does not change its big-$O$ behavior. This tradeoff --- reranking accuracy against reranking cost under a widened, fairness-motivated candidate pool --- is, in our reading, understated in reranking literature that evaluates cost only against a single-corpus, unwidened baseline.

\section{Chunking Strategies}
\label{sec:chunking}

Retrieval granularity is usually decided before any retrieval model is chosen: documents are split into fixed-size, overlapping chunks, and the chunk --- not the document --- becomes the unit that is embedded, indexed, and retrieved. Fixed-size chunking with a modest overlap remains the dominant practical choice because it is cheap, deterministic, and works passably across many document types, but its tradeoffs have increasingly been studied directly rather than assumed. Chunk-size sensitivity varies by domain and document structure, as shown for financial reports, where headers, tables, and dense terminology interact with the chosen chunk boundary in ways that affect downstream retrieval quality measurably \cite{jimenoyepes2024chunking}. ``Late chunking'' \cite{gunther2024latechunking} inverts the usual order of operations, encoding a long document with a long-context embedding model first and deriving chunk-level embeddings from token-level representations afterward, so that each chunk's embedding still reflects surrounding document context that a naively-chunked-then-embedded passage would lose. Comparative studies of chunking strategies more broadly \cite{merola2025chunking} suggest that no single strategy dominates across document types, which is precisely the situation a multi-source system is in by construction: a website's prose, a PDF's mixed prose-and-table layout, a video transcript's spoken, punctuation-light stream, and a source code file's line-oriented structure are different enough that a chunking strategy tuned for one is unlikely to be well-tuned for another. A practical implication we return to in Section~\ref{sec:casestudy} is that multi-source systems in practice tend to adopt a shared default chunker for prose-like sources and a distinct chunking unit specifically for whichever modality least resembles prose (most often code, by line, or transcript, by time/character band), rather than forcing one chunking function to serve every modality equally.

Chunk overlap deserves separate comment from chunk size, since the two are often conflated as one hyperparameter. Overlap exists specifically to prevent a fact that straddles a chunk boundary from being unrecoverable by either resulting chunk; the practical cost of overlap is redundant storage and redundant embedding computation proportional to the overlap fraction, which is a fixed, easily-budgeted cost, unlike chunk size, whose effect on retrieval quality is corpus- and query-dependent and harder to budget for in advance. A further wrinkle specific to multi-source systems is that overlap's benefit is not uniform across modalities: a transcript chunked by a time/character band, where segment boundaries already align with natural pauses in speech, needs overlap less than dense, boundary-agnostic prose extracted from a PDF's continuous text flow, which suggests --- though we are not aware of a study that measures this directly --- that a single global overlap parameter, like a single global chunk size, is a simplification a genuinely multi-modal chunking strategy would ultimately need to relax.

An alternative to fixed-size chunking altogether is to place chunk boundaries adaptively, at points where adjacent sentences' embeddings diverge beyond some threshold (``semantic chunking''), so that a chunk boundary tends to coincide with an actual topic shift rather than an arbitrary character count; late chunking \cite{gunther2024latechunking} is one instance of this broader adaptive-boundary family, and comparative evaluations \cite{merola2025chunking} generally find adaptive strategies competitive with, but not uniformly better than, well-tuned fixed-size chunking, at the cost of an extra embedding pass over candidate boundaries before the final chunk set is fixed. For a multi-source system, this cost tradeoff is itself modality-dependent: computing sentence-level embeddings to find adaptive boundaries is a comparatively small additional cost for prose, where sentence segmentation is already a solved subproblem, but is considerably less well-defined for source code, where a ``sentence'' has no clear analogue and syntax-aware boundaries (function or class scope) are already a more natural adaptive criterion than embedding divergence.

\section{Modality-Specific Retrieval Challenges}
\label{sec:modality}

\subsection{Web Documents}
Retrieval over crawled web content inherits two problems documents-at-rest do not have: content must first be extracted from a live, frequently JavaScript-rendered page rather than read from a static file, and a crawl frontier must be bounded (by domain, depth, or page count) to avoid retrieving and indexing arbitrary, possibly unbounded, linked content. Both concerns sit upstream of retrieval proper, but they directly shape what ends up retrievable: a page that fails to render, or that lies outside the crawl's domain or depth limit, is simply absent from the corpus, which is a silent failure mode specific to this modality that a static-document corpus does not share. A further, easily overlooked issue is near-duplicate content: templated navigation, footers, and boilerplate repeated across every page of a crawled site can dominate a naive chunk's lexical signature, diluting the very content that made the page worth crawling in the first place; systems that do not explicitly strip or down-weight repeated boilerplate risk retrieving several near-identical, low-information chunks from different pages of the same site rather than one genuinely relevant one.

\subsection{Video and Transcript QA}
Video question answering over transcripts introduces two distinguishing characteristics. First, the natural retrieval unit is time-aligned rather than purely textual: a useful citation to video evidence names a timestamp range, not merely a document title, which existing benchmarks such as TVQA \cite{lei2018tvqa} evaluate directly by requiring temporally localized evidence rather than whole-video labels. Second, when working from an automatically generated transcript rather than a human-authored one, retrieval-time text is noisier --- transcription errors, missing punctuation, imprecise sentence boundaries --- than the crawled or authored text a system's other sources might contain, which has consequences for how aggressively such a system should trust retrieval confidence computed the same way across all of its sources (Section~\ref{sec:confidence}). Frameworks such as LLoVi \cite{zhang2024llovi} address long-form video QA by first converting video to language (captions or transcripts) and then applying standard long-context language-model reasoning over that text, which is architecturally the same reduction --- treat video as retrievable text plus timestamps --- that a transcript-based multi-source system makes. A further practical tension specific to this modality is chunk-size choice under a timestamp-labeling constraint: a shorter transcript chunk gives a more precise citation timestamp but a weaker, more fragmentary retrieval signal, since spoken language rarely carries the same self-contained-sentence structure that written prose does; a longer chunk improves retrieval signal at the cost of a citation timestamp that only loosely localizes where in a multi-minute span the answer actually appears.

\subsection{Code-Repository Retrieval}
Code differs from prose in ways that affect both extraction and retrieval. Structurally, a repository is not one document but a directory tree of files with import/dependency relationships between them, and a chunk boundary that respects function or class scope is more useful than an arbitrary character cut, motivating line- or syntax-aware chunking rather than the character-budget chunking common for prose. Retrieval-wise, semantic code search was an early target for dense retrieval methods, as in the CodeSearchNet benchmark \cite{husain2019codesearchnet}, and has since extended into retrieval-augmented code completion that must decide not only what to retrieve but from \emph{where in the repository}: RepoCoder \cite{zhang2023repocoder} iterates retrieval and generation so that an initial (possibly imperfect) generation can itself be used to retrieve better subsequent context, and Repoformer \cite{wu2024repoformer} learns to selectively decide when repository-level retrieval is worth its cost at all versus when a model's own local context already suffices. A repository's own structure --- its README, its directory layout, its test and configuration conventions --- is itself a source of retrievable signal distinct from any individual file's content, a point we return to concretely in Section~\ref{sec:casestudy}. A further consideration, only partially addressed in the retrieval-augmented code literature surveyed above, is that a repository's files are not independent: an import or call-graph edge between two files is itself evidence of relatedness that pure text-similarity retrieval cannot see, so a chunk-level retriever operating on code in isolation from its dependency graph is discarding structural information a compiler or static analyzer already has for free.

\subsection{Multi-Document and Multi-Source Fusion}
Retrieving well from any single source is necessary but not sufficient once a query may be legitimately answered by evidence spread across several sources at once. End-to-end training of a reader and retriever jointly over multiple documents \cite{sachan2021multidoc} treats multi-document evidence aggregation as a first-class training objective rather than an artifact of independently ranking single documents and concatenating the top results. Reciprocal Rank Fusion \cite{cormack2009rrf}, discussed above as a reranking technique, is equally a multi-source fusion technique when its inputs are ranked lists retrieved independently \emph{per source} rather than fields of one corpus. A distinguishing practical requirement of multi-source (as opposed to merely multi-document) retrieval is source-level fairness: a user who explicitly selects several sources to query jointly has an expectation, often implicit, that each selected source was actually considered, which a pure top-$k$-by-score selection does not guarantee if one source's evidence is systematically weaker or sparser than another's --- a failure mode we discuss concretely, with a numeric example, in Section~\ref{sec:casestudy}. This requirement is orthogonal to, and can be layered on top of, whichever fusion technique is otherwise used: unweighted rank fusion, learned re-weighting, or a joint reader-retriever objective can all be combined with an explicit source-coverage constraint applied after scoring, since the constraint operates on \emph{which source produced each candidate}, a piece of provenance metadata that fusion techniques generally carry through unchanged but rarely act on directly themselves.

\section{Confidence, Calibration, and Abstention}
\label{sec:confidence}

A grounded QA system that always answers, regardless of retrieval quality, will eventually answer confidently from irrelevant or insufficient evidence. Selective prediction under distribution shift \cite{kamath2020selective} frames this as a decision the system should make explicitly --- answer, or abstain --- rather than as an emergent property of whatever the generator happens to produce. Whether a language model's own reported confidence can be trusted for this purpose is itself contested: \cite{kadavath2022knowwhattheyknow} finds that larger models can be reasonably well calibrated on some self-assessment tasks, but hallucination surveys \cite{ji2023hallucination} document broad and persistent failure modes in which fluent, confident-sounding generation is simply wrong, which argues for grounding the abstention decision in retrieval evidence rather than the generator's self-report wherever retrieval evidence is available. Retrieval augmentation has itself been shown to reduce hallucination directly in open-ended generation \cite{shuster2021retrieval}, which further supports gating on retrieval quality specifically, rather than only on downstream generation behavior after the fact.

Multi-source retrieval complicates calibration in a way single-corpus systems do not encounter: if different sources have systematically different retrieval-quality distributions --- an automatically transcribed video's noisier text against a crawled webpage's cleaner text, for instance --- then a single fixed confidence threshold, applied uniformly across all sources, will be miscalibrated for at least one of them. A system must either normalize confidence per source-type before thresholding, or make an explicit, stated exception for the source types where uniform thresholding is known to misfire; we discuss a concrete instance of the latter approach in Section~\ref{sec:casestudy}.

Abstention itself is not a binary design choice in practice: a system can refuse outright, answer with an attached caveat, or ask a clarifying follow-up question in place of either extreme. The middle option --- answering with a caveat rather than refusing --- trades a harder guarantee (no unsupported claims reach the user) for a softer one (unsupported claims are flagged, but still delivered), which may be the more defensible choice specifically for source types, such as noisy transcripts, where a strict refusal threshold would decline a disproportionate, and possibly majority, share of genuinely answerable questions. Which of these three responses a system should choose for a given confidence band is, to our knowledge, an underspecified design question in the literature we surveyed above; most selective-prediction work \cite{kamath2020selective} treats abstention as binary, leaving the graded, per-source-type policy question --- refuse, caveat, or clarify --- to individual system builders' judgment rather than to any established methodology.

\section{Case Study: SAGE}
\label{sec:casestudy}

SAGE (Semantic Analysis and Generation Engine) ingests four source types --- crawled multi-page websites, PDF documents, YouTube video transcripts, and GitHub code repositories --- through modality-specific extractors that each terminate in a common chunk representation, embeds all chunks with the same sentence encoder \cite{reimers2019sbert,wang2020minilm}, and answers questions against one, several, or all of a user's sources through a single retrieval, reranking, confidence-gating, and generation pipeline. We do not repeat SAGE's full architecture here; instead we use two of its design decisions as concrete instances of ideas introduced earlier in this survey, and note two further points briefly. Table~\ref{tab:casestudy} previews how each of SAGE's four modalities maps back onto the modality-specific challenges discussed in Section~\ref{sec:modality}.

\begin{table}[t]
\caption{SAGE's per-modality choices as instances of Section~\ref{sec:modality}'s challenges}
\label{tab:casestudy}
\centering
\scriptsize
\begin{tabular}{@{}>{\raggedright\arraybackslash}p{1.3cm}>{\raggedright\arraybackslash}p{3.0cm}>{\raggedright\arraybackslash}p{3.4cm}@{}}
\toprule
Modality & Survey challenge (\S\ref{sec:modality}) & SAGE's instantiation \\
\midrule
Website & bounded crawl frontier, boilerplate dilution & same-domain BFS crawl; per-page overview excerpt reduces reliance on any one boilerplate-heavy chunk \\
Video & timestamp-precision/signal tradeoff, noisier ASR text & fixed char-band transcript chunks with timestamp metadata; confidence exception for video (\S\ref{sec:confidence}) \\
Code & structure beyond individual files, line- not character-scale chunks & line-range chunking; filesystem-introspective overview chunk surfaces directory structure and conventions directly \\
Multi-source & source-level fairness under score-only ranking & diversity floor reserving one slot per explicitly combined source before score-filling remaining slots \\
\bottomrule
\end{tabular}
\end{table}

\textbf{Overview chunks as a chunking-strategy instance.} Section~\ref{sec:chunking} noted that no single chunking strategy dominates across the modalities a multi-source system ingests, and that broad, non-specific queries pose a particular retrieval problem: a question like ``what is this about'' has poor lexical and semantic overlap with any single content chunk, because the answer lives in a source's overall structure rather than any one passage. SAGE addresses this not by choosing a cleverer chunk size, but by adding a synthetic, deterministic \emph{overview chunk} to every source at ingestion time --- built from the source's own statistics and excerpts (crawled page titles, transcript timestamps, or, for a GitHub repository, a live-introspected on-disk directory tree and detected project conventions), with no additional generative model call. This is a chunking-strategy decision in the sense of Section~\ref{sec:chunking} --- it is, in effect, a fifth chunk added to every source specifically to change what is retrievable --- rather than a generation-time summarization feature, and it requires no labeled training data or model fine-tuning to implement per new modality.

\textbf{The diversity floor as a multi-source-fusion instance.} Section~\ref{sec:modality} identified source-level fairness as a requirement specific to multi-source (not merely multi-document) retrieval: a query that explicitly names several sources should not silently omit one of them from its evidence because that source's single best-matching chunk scored lower than another source's several good chunks. SAGE encountered this failure mode directly during its own combined-source testing --- a two-source query citing only one of the two sources it named, despite both containing genuinely relevant content --- and addressed it with a diversity floor: before filling a fixed number of citation slots by pure reranked score, each explicitly combined source's single best-scoring chunk is reserved first, and only the remaining slots are filled by score across all sources. Concretely, given five candidate chunks with blended scores 0.7326, 0.7064, 0.6758, 0.6652, and 0.4730 drawn two from one source and three from another, a pure score cutoff at $k{=}2$ happens to keep one chunk per source in this instance, but only because the second source's best chunk narrowly outscored the first source's second- and third-best; a small perturbation of that single score (a lexical-overlap difference of 0.3, in this case) would drop the second source out of the answer entirely under pure score ranking, while the reservation step prevents this unconditionally rather than by chance. We consider this a useful concrete data point for the general claim in Section~\ref{sec:modality}: source-level fairness failures in multi-source retrieval are not obviously wrong when reranking is examined in isolation, and tend to surface only once a system is tested against genuinely combined, multi-source queries directly.

\textbf{Confidence gating with a stated per-modality exception.} Consistent with the calibration concern raised in Section~\ref{sec:confidence}, SAGE computes a single retrieval-derived confidence score from its top-ranked and top-three blended scores plus a small evidence-count bonus, and refuses to answer in its grounded mode below a fixed threshold — except for chunks originating from video sources, which are exempted from hard refusal (receiving a lower-confidence warning instead), on the stated reasoning that automatically transcribed text is measurably noisier than the system's other, cleaner-text sources. This is the explicit-per-source-exception strategy anticipated in Section~\ref{sec:confidence}, applied to exactly one modality rather than derived from a general per-source calibration procedure.

\textbf{Evidence, not a benchmark.} We report SAGE's own test suite (95 test functions across 12 files, covering chunking, URL normalization, reranker and confidence formula sanity checks, per-modality overview-chunk generation, and an end-to-end API integration suite) as evidence that the mechanisms above are implemented and exercised, not as a quantitative retrieval-quality or answer-quality result; SAGE does not report a benchmark score against a labeled dataset, and we do not construct one for it here, consistent with Section~\ref{sec:challenges}'s point that a combined-modality benchmark of the kind that would support such a score does not, to our knowledge, yet exist.

\textbf{Minor points.} Two aspects of SAGE that are not this survey's focus deserve one sentence each for completeness. SAGE mitigates server-side request forgery in its website crawler by validating resolved IP addresses against private/reserved ranges on every followed redirect, not only the entry URL \cite{owasp2021ssrf}, and mitigates indirect prompt injection from untrusted retrieved web content by instructing its generator explicitly to treat retrieved context as quoted data rather than instructions \cite{greshake2023promptinjection}. SAGE is also comparable, at a product level, to consumer ``chat with your sources'' tools such as NotebookLM \cite{notebooklm2023} and to general-purpose composability frameworks such as LangChain \cite{langchain} and LlamaIndex \cite{llamaindex}, but occupies a different point in the design space from either --- a specific, opinionated multi-source pipeline rather than a polished single-user product or a generic library --- and we do not attempt a feature-by-feature comparison here.

\section{Open Challenges}
\label{sec:challenges}

Several problems recur across the systems and papers surveyed above without a settled answer. \emph{Cross-modal fusion at the scoring level} remains largely heuristic: blended semantic-lexical scores, rank fusion, and learned combination methods (Section~\ref{sec:reranking}) each make different, usually untested, assumptions about whether a score computed from one modality's chunk is comparable in magnitude to a score computed from another's; SAGE's diversity floor sidesteps this by guaranteeing representation rather than by claiming cross-modality score comparability, which trades a stronger guarantee (fairness) for a weaker one (true relevance ranking across sources). \emph{Confidence calibration across heterogeneous corpora} is similarly unresolved in general: per-source-type exceptions, as used for video transcripts in Section~\ref{sec:casestudy}, are a practical patch rather than a calibrated solution, and we are not aware of published work that calibrates a single retrieval-confidence score jointly across genuinely heterogeneous modalities with a labeled dataset spanning all of them. \emph{Incremental re-indexing} is underexplored specifically for multi-source systems: chunking strategies (Section~\ref{sec:chunking}) are usually described as one-time operations, but a source that changes over time (a recrawled website, an updated repository) raises the question of which existing chunks should be preserved and which re-embedded, a cost/correctness tradeoff that content-addressed chunk identifiers can make explicit but do not resolve on their own. \emph{Evaluation} for multi-source, multi-modal RAG lags behind single-corpus RAG evaluation: most retrieval and QA benchmarks (including several surveyed in Section~\ref{sec:modality}) are modality-specific, and we are not aware of a widely adopted benchmark that jointly measures retrieval quality, citation correctness, and confidence calibration across web, document, video, and code sources combined in a single query, which would be necessary to compare systems like SAGE against one another on genuinely equal terms rather than on architecture description alone. Standard retrieval metrics compound this gap rather than closing it: Recall@$k$, MRR, and nDCG are all defined against a single, fixed relevance judgment per query, which presumes a labeled query-document(s) pair drawn from one corpus; extending them to a combined-source query requires deciding, before any metric can even be computed, whether relevance judgments should be pooled across sources or required per source (to capture the fairness property discussed in Sections~\ref{sec:modality} and \ref{sec:casestudy}), and we are not aware of a standard convention for making that choice.

\emph{Fine-tuning against a multi-modal negative distribution}, raised in Section~\ref{sec:foundations}, is a training-side gap rather than an architectural one: nearly every open-source sentence encoder in wide use, including the one our case study uses, was trained on question-passage pairs drawn from a single kind of corpus, and we are not aware of published work that deliberately constructs hard negatives spanning prose, transcript, and code together to fine-tune a single encoder for genuinely mixed-modality retrieval, as opposed to fine-tuning separate encoders per modality and reconciling their outputs post hoc.

\emph{Dependency-aware retrieval for structured sources}, raised in Section~\ref{sec:modality}'s discussion of code, generalizes beyond code: a website's internal link graph and a PDF's cross-references are analogous structural signals that most chunk-level dense retrievers discard identically to how they discard a repository's import graph, suggesting that the code-specific graph-aware retrieval methods surveyed in Section~\ref{sec:modality} may be a more general opportunity for structured, multi-source retrieval than their code-specific framing in the literature suggests.

\section{Conclusion}
\label{sec:conclusion}

Grounded question answering over heterogeneous, multi-source knowledgebases draws on the same core techniques as single-corpus RAG --- dense retrieval, hybrid reranking, chunking, and confidence-gated generation --- but requires each of them to be reconsidered once ``the corpus'' is actually several structurally different corpora combined at query time. Chunking strategies must accommodate prose, transcripts, and code without assuming one size fits all; reranking must be paired with an explicit source-fairness mechanism once queries can name several sources at once; and confidence estimation must either normalize across source-type-specific retrieval-quality distributions or state its exceptions plainly. We used SAGE's overview-chunk mechanism and diversity floor as concrete, checkable instances of the second and third of these adjustments, not as a departure from the broader literature but as one system's specific resolution of tradeoffs that literature leaves open. Cross-modal score comparability, joint confidence calibration, incremental re-indexing, and combined-modality evaluation remain, in our reading, the open problems most in need of dedicated attention as multi-source RAG systems move from research prototypes toward more general use.

\begin{thebibliography}{00}

\bibitem{lewis2020rag} P. Lewis, E. Perez, A. Piktus, F. Petroni, V. Karpukhin, N. Goyal, H. K\"{u}ttler, M. Lewis, W. Yih, T. Rockt\"{a}schel, S. Riedel, and D. Kiela, ``Retrieval-augmented generation for knowledge-intensive NLP tasks,'' in \emph{Advances in Neural Information Processing Systems (NeurIPS)}, 2020. arXiv:2005.11401.

\bibitem{karpukhin2020dpr} V. Karpukhin, B. O\u{g}uz, S. Min, P. Lewis, L. Wu, S. Edunov, D. Chen, and W. Yih, ``Dense passage retrieval for open-domain question answering,'' in \emph{Proc. EMNLP}, 2020, pp. 6769--6781. arXiv:2004.04906.

\bibitem{izacard2021fid} G. Izacard and E. Grave, ``Leveraging passage retrieval with generative models for open domain question answering,'' in \emph{Proc. EACL}, 2021. arXiv:2007.01282.

\bibitem{gao2023ragsurvey} Y. Gao, Y. Xiong, X. Gao, K. Jia, J. Pan, Y. Bi, Y. Dai, J. Sun, M. Wang, and H. Wang, ``Retrieval-augmented generation for large language models: A survey,'' 2023. arXiv:2312.10997.

\bibitem{reimers2019sbert} N. Reimers and I. Gurevych, ``Sentence-BERT: Sentence embeddings using Siamese BERT-networks,'' in \emph{Proc. EMNLP-IJCNLP}, 2019. arXiv:1908.10084.

\bibitem{wang2020minilm} W. Wang, F. Wei, L. Dong, H. Bao, N. Yang, and M. Zhou, ``MiniLM: Deep self-attention distillation for task-agnostic compression of pre-trained transformers,'' in \emph{NeurIPS}, 2020. arXiv:2002.10957.

\bibitem{wang2020minilmv2} W. Wang, H. Bao, S. Huang, L. Dong, and F. Wei, ``MiniLMv2: Multi-head self-attention relation distillation for compressing pretrained transformers,'' 2020. arXiv:2012.15828.

\bibitem{minilm_modelcard} Sentence-Transformers, ``all-MiniLM-L6-v2 model card,'' Hugging Face. [Online]. Available: https://huggingface.co/sentence-transformers/all-MiniLM-L6-v2

\bibitem{robertson2009bm25} S. Robertson and H. Zaragoza, ``The probabilistic relevance framework: BM25 and beyond,'' \emph{Foundations and Trends in Information Retrieval}, vol. 3, no. 4, pp. 333--389, 2009.

\bibitem{nogueira2019bert} R. Nogueira and K. Cho, ``Passage re-ranking with BERT,'' 2019. arXiv:1901.04085.

\bibitem{khattab2020colbert} O. Khattab and M. Zaharia, ``ColBERT: Efficient and effective passage search via contextualized late interaction over BERT,'' in \emph{Proc. 43rd Int. ACM SIGIR Conf.}, 2020. arXiv:2004.12832.

\bibitem{gao2020clear} L. Gao, Z. Dai, and J. Callan, ``Complement lexical retrieval model with semantic residual embeddings,'' 2020. arXiv:2004.13969.

\bibitem{cormack2009rrf} G. V. Cormack, C. L. A. Clarke, and S. B\"{u}ttcher, ``Reciprocal rank fusion outperforms Condorcet and individual rank learning methods,'' in \emph{Proc. 32nd Int. ACM SIGIR Conf.}, 2009, pp. 758--759.

\bibitem{jimenoyepes2024chunking} A. Jimeno-Yepes, Y. You, J. Milczek, S. Laverde, and R. Li, ``Financial report chunking for effective retrieval augmented generation,'' 2024. arXiv:2402.05131.

\bibitem{gunther2024latechunking} M. G\"{u}nther, I. Mohr, D. J. Williams, B. Wang, and H. Xiao, ``Late chunking: Contextual chunk embeddings using long-context embedding models,'' 2024. arXiv:2409.04701.

\bibitem{merola2025chunking} C. Merola and J. Singh, ``Reconstructing context: Evaluating advanced chunking strategies for retrieval-augmented generation,'' in \emph{2nd Workshop on Knowledge-Enhanced Information Retrieval, ECIR}, 2025. arXiv:2504.19754.

\bibitem{lei2018tvqa} J. Lei, L. Yu, M. Bansal, and T. L. Berg, ``TVQA: Localized, compositional video question answering,'' in \emph{Proc. EMNLP}, 2018. arXiv:1809.01696.

\bibitem{zhang2024llovi} C. Zhang, T. Lu, M. M. Islam, Z. Wang, S. Yu, M. Bansal, and G. Bertasius, ``A simple LLM framework for long-range video question-answering,'' in \emph{Proc. EMNLP}, 2024. arXiv:2312.17235.

\bibitem{husain2019codesearchnet} H. Husain, H.-H. Wu, T. Gazit, M. Allamanis, and M. Brockschmidt, ``CodeSearchNet challenge: Evaluating the state of semantic code search,'' 2019. arXiv:1909.09436.

\bibitem{zhang2023repocoder} F. Zhang, B. Chen, Y. Zhang, J. Keung, J. Liu, D. Zan, Y. Mao, J.-G. Lou, and W. Chen, ``RepoCoder: Repository-level code completion through iterative retrieval and generation,'' in \emph{Proc. EMNLP}, 2023. arXiv:2303.12570.

\bibitem{wu2024repoformer} D. Wu, W. U. Ahmad, D. Zhang, M. K. Ramanathan, and X. Ma, ``Repoformer: Selective retrieval for repository-level code completion,'' in \emph{Proc. ICML}, 2024. arXiv:2403.10059.

\bibitem{sachan2021multidoc} D. S. Sachan, S. Reddy, W. Hamilton, C. Dyer, and D. Yogatama, ``End-to-end training of multi-document reader and retriever for open-domain question answering,'' in \emph{NeurIPS}, 2021. arXiv:2106.05346.

\bibitem{kamath2020selective} A. Kamath, R. Jia, and P. Liang, ``Selective question answering under domain shift,'' in \emph{Proc. ACL}, 2020. arXiv:2006.09462.

\bibitem{kadavath2022knowwhattheyknow} S. Kadavath, T. Conerly, A. Askell, T. Henighan, D. Drain, E. Perez, N. Schiefer, Z. Hatfield-Dodds, N. DasSarma, E. Tran-Johnson, S. Johnston, S. El-Showk, A. Jones, N. Elhage, T. Hume, A. Chen, Y. Bai, S. Bowman, S. Fort, D. Ganguli, D. Hernandez, J. Jacobson, J. Kernion, S. Kravec, L. Lovitt, K. Ndousse, C. Olsson, S. Ringer, D. Amodei, T. Brown, J. Clark, N. Joseph, B. Mann, S. McCandlish, C. Olah, and J. Kaplan, ``Language models (mostly) know what they know,'' 2022. arXiv:2207.05221.

\bibitem{ji2023hallucination} Z. Ji, N. Lee, R. Frieske, T. Yu, D. Su, Y. Xu, E. Ishii, Y. J. Bang, A. Madotto, and P. Fung, ``Survey of hallucination in natural language generation,'' \emph{ACM Computing Surveys}, vol. 55, no. 12, art. 248, 2023.

\bibitem{shuster2021retrieval} K. Shuster, S. Poff, M. Chen, D. Kiela, and J. Weston, ``Retrieval augmentation reduces hallucination in conversation,'' in \emph{Findings of EMNLP}, 2021, pp. 3784--3803. arXiv:2104.07567.

\bibitem{huang2013dssm} P.-S. Huang, X. He, J. Gao, L. Deng, A. Acero, and L. Heck, ``Learning deep structured semantic models for web search using clickthrough data,'' in \emph{Proc. 22nd ACM Int. Conf. on Information and Knowledge Management (CIKM)}, 2013, pp. 2333--2338.

\bibitem{owasp2021ssrf} OWASP Foundation, ``A10:2021 -- Server-Side Request Forgery (SSRF),'' OWASP Top 10:2021. [Online]. Available: https://owasp.org/Top10/2021/A10\_2021-Server-Side\_Request\_Forgery\_(SSRF)/

\bibitem{greshake2023promptinjection} K. Greshake, S. Abdelnabi, S. Mishra, C. Endres, T. Holz, and M. Fritz, ``Not what you've signed up for: Compromising real-world LLM-integrated applications with indirect prompt injection,'' in \emph{Proc. 16th ACM Workshop on Artificial Intelligence and Security (AISec)}, 2023. arXiv:2302.12173.

\bibitem{notebooklm2023} Google, ``NotebookLM: How to try Google's experimental AI-first notebook,'' Google Blog, 2023. [Online]. Available: https://blog.google/technology/ai/notebooklm-google-ai/

\bibitem{langchain} LangChain, Inc., ``LangChain documentation.'' [Online]. Available: https://docs.langchain.com/

\bibitem{llamaindex} LlamaIndex, ``LlamaIndex documentation.'' [Online]. Available: https://developers.llamaindex.ai/python/framework/

\end{thebibliography}

\end{document}
